\documentclass[11pt,a4paper]{article}

\usepackage[margin=1in]{geometry}
\usepackage{amsmath,amssymb,amsthm}
\usepackage{microtype}
\usepackage{booktabs}
\usepackage{tabularx}
\usepackage{enumitem}
\usepackage{hyperref}
\usepackage{xcolor}
\usepackage{listings}
\usepackage{cite}

\hypersetup{
  colorlinks=true,
  linkcolor=blue!50!black,
  citecolor=blue!50!black,
  urlcolor=blue!50!black,
}

\newtheorem{theorem}{Theorem}
\newtheorem{lemma}[theorem]{Lemma}
\newtheorem{proposition}[theorem]{Proposition}
\theoremstyle{definition}
\newtheorem{definition}[theorem]{Definition}
\newtheorem{construction}[theorem]{Construction}
\theoremstyle{remark}
\newtheorem{remark}[theorem]{Remark}

\newcommand{\zap}{\textsf{ZAP}\xspace}
\newcommand{\zapmcp}{\textsf{ZAP-MCP}\xspace}
\newcommand{\zapaa}{\textsf{ZAP-A2A}\xspace}
\newcommand{\mcp}{\textsf{MCP}\xspace}
\newcommand{\aatwo}{\textsf{A2A}\xspace}
\newcommand{\rns}{\textsf{ZAP-RNS}\xspace}
\newcommand{\Adv}{\mathsf{Adv}}
\newcommand{\Exp}{\mathsf{Exp}}
\newcommand{\AEAD}{\textsf{AE}}
\newcommand{\KEM}{\textsf{KEM}}
\newcommand{\Sign}{\textsf{Sig}}
\newcommand{\negl}{\mathsf{negl}}

\usepackage{xspace}

\title{\textbf{\zapmcp{} and \zapaa{}: Mutually Authenticated, Non-Repudiable
Transport for the \mcp{} and Agent2Agent Protocols}}

\author{ZAP Protocol Authors\\
\href{https://zap-proto.io}{\texttt{zap-proto.io}}}

\date{2026}

\begin{document}
\maketitle

\begin{abstract}
Autonomous and semi-autonomous AI agents have moved from research demos to
infrastructure-critical systems: they execute trades, retrieve patient
records, and modify production code. The two protocols that knit these
agents together --- \emph{Model Context Protocol} (\mcp{}) for
agent-to-tool calls~\cite{mcp2024} and \emph{Agent2Agent} (\aatwo{}) for
agent-to-agent collaboration~\cite{a2a2025} --- inherit their security
posture from the legacy web stack: TLS for confidentiality, API keys or
JWT bearers for authentication, and HTTP/SSE or HTTP/JSON framing. We
argue that this stack is structurally inadequate for an agent
fabric. Bearer tokens are bearer-authenticated, not message-authenticated;
JSON-RPC payloads carry no identity; and intermediate agent gateways or
tool brokers are forced to terminate TLS in order to route, breaking
end-to-end confidentiality and provenance. We present \zapmcp{} and
\zapaa{}, two compositions of the \zap{} post-quantum authenticated
transport with the \mcp{} and \aatwo{} application protocols. Each
agent-to-tool or agent-to-agent message is wrapped in a \zap{} frame,
keyed to a Module-Lattice KEM public key, signed under a per-agent
ML-DSA-65 (or hybrid Ed25519+ML-DSA-65) signing key, and bound to a
session-scoped nonce ladder. We give security definitions, prove
non-repudiation under the EUF-CMA assumption for ML-DSA, prove
agent-impersonation infeasibility under combined KEM IND-CCA and signature
EUF-CMA assumptions, and bound replay advantage by the nonce-collision
probability. We also analyse wire overhead and show that \zap{}-wrapped
\mcp{}/\aatwo{} traffic is shorter on the wire than HTTP+SSE+JWT
baselines for typical tool-call sizes, while delivering audit-grade
provenance suitable for HIPAA, SOC~2, and GDPR Article~30 records of
processing.
\end{abstract}

\section{Introduction}

\subsection{Agents are now critical infrastructure}

A 2023 LLM-integrated chatbot is recoverable; a 2026 agent that has
funded your accounts, called your APIs, and rewritten your production
code is not. Agents are increasingly given tools that perform irreversible
actions: order routing, database mutations, identity issuance, code
deployment, fund transfers, prescription writing. The protocols that
carry these actions --- \mcp{} for tool invocation and \aatwo{} for
inter-agent collaboration --- are the new universal joints of automated
labor.

These protocols were specified for fast adoption, not for long-horizon
auditability or post-quantum survivability. \mcp{} runs JSON-RPC 2.0 over
stdio (no auth), HTTP+SSE (auth via TLS + an opaque ``API key'' header or
OAuth bearer token), or WebSocket (likewise)~\cite{mcp2024,jsonrpc2}.
\aatwo{} runs structured tasks and artifacts over HTTP/JSON, with agent
identification expressed via ``agent cards'' typically signed only by the
TLS certificate of the host~\cite{a2a2025}. In both cases, the security
boundary is the TLS connection, not the message.

\subsection{Why the TLS-and-bearer stack is inadequate}

We identify five structural gaps:

\begin{enumerate}[leftmargin=2em,itemsep=0.25em]
\item \textbf{Tool impersonation.} A bearer token attests that the
\emph{client} has presented credentials to \emph{some} server reached at
a URL. It does not attest that the server is the intended tool. URL
phishing or DNS substitution succeeds whenever the agent has not
out-of-band pinned the tool's certificate.
\item \textbf{Message non-repudiation.} TLS protects messages
\emph{in flight}, not afterwards. A logged JSON-RPC envelope cannot
later be verified to have actually been emitted by a given agent: TLS
authentication is symmetric session-keyed and unprovable to a third
party. Disputes over ``did agent $A$ really request to wire \$10\,M?''
have no cryptographic answer.
\item \textbf{Cross-session replay.} A bearer-token-authenticated request
captured at runtime can be replayed by an attacker with the same token,
because the request itself is not bound to a fresh session secret.
\item \textbf{Quantum-vulnerable transcripts.} Long-lived agent logs
(medical, legal, financial) recorded today and decryptable in 2035 by an
attacker holding a CRQC are out of policy under
``harvest-now-decrypt-later''~\cite{harvest_now} for any organisation
whose data retention exceeds the quantum horizon.
\item \textbf{Broker-opaque routing.} Gateways that route between agent
clouds must today terminate TLS in order to read \mcp{}/\aatwo{}
message envelopes; this destroys end-to-end secrecy and forces gateways
into the trusted compute base.
\end{enumerate}

\subsection{Thesis}

Wrapping \mcp{} and \aatwo{} in \zap{} gives every agent-to-tool and
agent-to-agent message: post-quantum confidentiality, mutual
cryptographic identity, non-repudiation, and provenance --- all at lower
wire overhead than the current HTTP+SSE / HTTP+JWT defaults. Concretely:

\begin{itemize}[leftmargin=2em,itemsep=0.25em]
\item Each \mcp{} JSON-RPC request and response, and each \aatwo{} task
or artifact, is encapsulated as a single \zap{} frame.
\item Each frame carries an ML-DSA-65 signature of the request hash by
the agent's long-term signing key (or a hybrid Ed25519+ML-DSA-65 pair
during transition).
\item Each frame is encrypted to the recipient's ML-KEM-768 public key
via an X-Wing-style hybrid construction~\cite{kyber2018,nistmlkem}.
\item The session is initiated by a \zap{} 1-RTT handshake whose
transcript is signed by both parties; thereafter, frames carry a
sequence-bound nonce.
\item Agent identity is resolved through \rns{}~\cite{zaprns}: a name
maps to an ML-KEM public key and a signing public key, replacing the
``URL plus host certificate'' identity model.
\end{itemize}

The remainder of this paper formalises the construction
(Sections~\ref{sec:threat}--\ref{sec:zap-a2a}), states and proves
non-repudiation, replay resistance, and impersonation resistance
(Section~\ref{sec:security}), measures wire overhead
(Section~\ref{sec:perf}), and discusses compliance, deployment, and
gap items (Sections~\ref{sec:compliance}--\ref{sec:gaps}).

\section{Background}

\subsection{Model Context Protocol}

\mcp{}~\cite{mcp2024} is a JSON-RPC 2.0~\cite{jsonrpc2} request/response
protocol between an \emph{agent} (client, typically driven by an LLM)
and a \emph{tool server}. The agent invokes methods such as
\texttt{tools/list}, \texttt{tools/call}, \texttt{resources/read}, and
\texttt{prompts/get}; the tool server returns structured results.
Three transports are specified: stdio (in-process), HTTP+SSE (request
by HTTP POST, asynchronous server messages by Server-Sent
Events~\cite{rfc6455}), and WebSocket. Authentication is undefined at
the protocol level; in practice, deployments inject an
\texttt{Authorization} header containing an OAuth 2.0 bearer
token~\cite{rfc6749,rfc7519}.

\subsection{Agent2Agent}

\aatwo{}~\cite{a2a2025} is a higher-level protocol over HTTP/JSON.
Agents publish \emph{agent cards} describing their capabilities; clients
invoke long-running \emph{tasks}, which produce
\emph{artifacts} consumed by other agents or terminating in a final
result. \aatwo{} authentication is HTTP-layer, typically TLS server
authentication plus a bearer JWT.

\subsection{The \zap{} transport}

\zap{}~\cite{zaptransport} is a binary, stateful, post-quantum
authenticated transport. A \zap{} session is established by a 1-RTT
handshake derived from the SIGMA~\cite{sigma} family, parameterised by
a hybrid KEM (ML-KEM-768 with X25519~\cite{kyber2018,nistmlkem}) and a
hybrid signature (Ed25519 with ML-DSA-65~\cite{nistmldsa,ed25519}).
Frames carry a 64-bit sequence number, a 96-bit nonce derived from
HKDF~\cite{krawczyk2010hkdf} of the session key, and an authenticated
ciphertext under AES-256-GCM or ChaCha20-Poly1305. Optionally a
\emph{signed-payload} flag adds a per-frame ML-DSA-65 signature over
the plaintext payload hash, providing third-party verifiable
non-repudiation.

\subsection{\rns{}: identity-bound naming}

\rns{}~\cite{zaprns} resolves a human-readable agent or tool name to a
record containing (i) an ML-KEM public key, (ii) an Ed25519 public key,
(iii) an ML-DSA-65 public key, (iv) a vector of capability claims, and
(v) a signature chain rooted at a publishing authority. Resolution
returns a self-authenticating record; identity is no longer expressed
as ``hostname + a certificate someone in the chain trusts''.

\section{Threat Model}\label{sec:threat}

\subsection{Trust assumptions}

\textbf{Trusted:} the agent's long-term signing key and the tool's
long-term decapsulation key (each held in an HSM or equivalent
isolated enclave); the \rns{} record cache after first authenticated
fetch; the cryptographic primitives ML-DSA-65, ML-KEM-768, AES-GCM,
SHA-3, HKDF.

\textbf{Untrusted:} the network; any intermediate ``agent gateway''
performing routing or load balancing; the JSON-RPC transcript on disk
(must be third-party verifiable); the DNS resolver (because identity
binds to KEM PK, not DNS).

\subsection{Adversary capabilities}

We assume an active network adversary $\mathcal{A}$ who:
\begin{itemize}[leftmargin=2em,itemsep=0.25em]
\item observes, drops, reorders, modifies, and injects all
network traffic;
\item operates a polynomial number of legitimate-looking tool servers
and agents and can cause them to engage in handshakes;
\item may corrupt up to all but one of the parties in the system,
provided we explicitly declare the surviving party honest;
\item is bounded to probabilistic polynomial time, including
polynomially many quantum oracle queries (the assumption used by
ML-KEM-768 and ML-DSA-65 security claims~\cite{nistmlkem,nistmldsa});
\item may store transcripts indefinitely (harvest-now-decrypt-later).
\end{itemize}

\subsection{Concrete agent-fabric attacks}

We consider the following attacks of practical importance, and the
section that addresses each:

\begin{enumerate}[leftmargin=2em,itemsep=0.25em]
\item \emph{Tool impersonation.} $\mathcal{A}$ stands up a malicious
tool server claiming to be a legitimate API.
\textbf{Addressed by} \rns{}-bound mutual authentication
(\S\ref{sec:zap-mcp}).
\item \emph{Indirect prompt injection via tool
output}~\cite{prompt_injection_2023}. Out of scope cryptographically
(prompt injection lives at the model layer); however, \zap{} provides
the prerequisite that tool output is cryptographically attributable to
a named tool, so injection-induced damage is forensically
reconstructible.
\item \emph{Cross-agent replay.} $\mathcal{A}$ captures a tool call from
agent $A$ and replays it as if from $A$, on a different session.
\textbf{Addressed by} session-bound nonces
(Lemma~\ref{lem:replay} in \cite{proof}).
\item \emph{Long-term provenance loss.} Six months later, agent $A$
denies having issued a transaction. \textbf{Addressed by} ML-DSA-65
signed frames (Theorem~\ref{thm:nonrep}).
\item \emph{Future-quantum decryption} of medical/legal transcripts.
\textbf{Addressed by} hybrid ML-KEM-768 + X25519 frame keying.
\item \emph{Broker-introduced injection.} A routing gateway modifies
\mcp{} request fields. \textbf{Addressed by} per-frame AEAD: a routing
gateway cannot decrypt or modify without invalidating the frame tag
and signature.
\end{enumerate}

\section{\zapmcp{} Construction}\label{sec:zap-mcp}

\subsection{Wire mapping}

A \mcp{} JSON-RPC request is a JSON object
$\rho = \{\texttt{jsonrpc},\texttt{id},\texttt{method},\texttt{params}\}$.
Define the canonical encoding
\[
\hat{\rho} = \mathrm{CBOR}(\rho)
\]
using deterministic CBOR~\cite{cose}. Let
$\sigma = \Sign.\mathrm{Sign}(sk_A, h(\hat{\rho}\;\Vert\;\mathit{ctx}))$
where $h$ is SHA3-256, $sk_A$ is agent $A$'s ML-DSA-65 signing key, and
$\mathit{ctx}$ is a binding context string of the form
\[
\mathit{ctx} = \mathtt{"zap-mcp/v1"} \,\Vert\, \mathit{sid} \,\Vert\, \mathit{seq}
\]
where $\mathit{sid}$ is the 32-byte session identifier produced by the
\zap{} handshake and $\mathit{seq}$ is the 64-bit frame sequence number.

The wrapped request is the \zap{} frame
\[
F = \zap.\mathrm{Seal}(K_{AB},\, \mathit{nonce}(\mathit{sid},\mathit{seq}),\,
\hat{\rho}\;\Vert\;\sigma,\, \mathit{ad})
\]
where $K_{AB}$ is the per-direction AEAD key derived from the handshake,
and $\mathit{ad} = \mathtt{"zap-mcp/v1/req"} \,\Vert\, \mathit{seq}$.

The corresponding \mcp{} response from tool $T$ is wrapped symmetrically
under $T$'s signing key, with $\mathit{ctx}$ extended to bind to the
request hash, ensuring response-to-request linkage.

\subsection{Mutual identification}

The agent locates the tool by name $n$ via \rns{}: it retrieves a record
$R_T$ binding $n$ to a tuple
$(\mathit{kem\_pk}_T, \mathit{sig\_pk}_T, \mathit{caps}_T)$ with a
publisher signature. The agent verifies the publisher chain, then
initiates a \zap{} handshake against the tool's network endpoint using
$\mathit{kem\_pk}_T$. The handshake transcript is signed by both
parties.

\begin{remark}
Crucially, the agent never trusts a hostname. If DNS misroutes the
agent to an attacker's endpoint, the handshake fails because the
attacker does not hold the secret key for $\mathit{kem\_pk}_T$.
This eliminates the entire class of tool-URL phishing attacks.
\end{remark}

\subsection{Capability claims}

\rns{} records carry capability claims, e.g.\ ``this tool is permitted
to list orders but not place them''. The handshake transcript binds the
hash of the claims set; the tool must therefore present a record whose
claims include the methods the agent will call. The agent can refuse
\texttt{tools/call} for methods absent from the claim vector before any
network round trip.

\subsection{Pseudocode (request path)}

\begin{lstlisting}[basicstyle=\footnotesize\ttfamily,frame=single,
                   xleftmargin=1em,framesep=0.5em]
function zapmcp_call(agent_sk, tool_name, method, params):
    R_T  := rns_resolve(tool_name)
    verify_record_chain(R_T)
    require method in R_T.caps

    sess := zap_handshake(R_T.kem_pk, agent_sk)  // 1-RTT
    rho  := { jsonrpc: "2.0", id: fresh_id(),
              method: method, params: params }
    rho_hat := cbor(rho)
    ctx  := "zap-mcp/v1" || sess.id || sess.next_seq()
    sig  := mldsa65_sign(agent_sk, sha3_256(rho_hat || ctx))
    F    := zap_seal(sess, rho_hat || sig)
    send(F)
    F'   := recv()
    payload, sig' := zap_open(sess, F')
    require mldsa65_verify(R_T.sig_pk,
                           sha3_256(payload || resp_ctx), sig')
    return cbor_decode(payload)
\end{lstlisting}

\section{\zapaa{} Construction}\label{sec:zap-a2a}

\aatwo{} is task-shaped, not call-shaped: an agent submits a task
description, the remote agent acknowledges and emits a sequence of
artifacts and progress events, and eventually a final result. We map
each top-level \aatwo{} message --- task envelope, artifact, status
event, terminal result --- to one signed \zap{} frame, with the
following augmentations.

\subsection{Signed agent cards}

In \aatwo{} an \emph{agent card} declares an agent's identity,
capabilities, and endpoints. Under \zapaa{}, an agent card is a
canonical CBOR object signed by the agent's ML-DSA-65 key and published
through \rns{}. A consuming agent verifies the card off the wire; the
card is itself a frame-shaped, third-party-verifiable artifact.

\subsection{Artifact provenance}

Each artifact emitted during a task carries a frame signature by the
producing agent. Because frames bind $\mathit{sid}$ (session id) and
the parent task identifier, an artifact is independently verifiable as
``produced by agent $B$, in the context of task $\tau$, in session
$\mathit{sid}$, at sequence $\mathit{seq}$''. A verifier presented with
only the artifact and the public keys of $B$ can confirm authenticity
without access to the original session.

\subsection{Multi-hop tasks}

When agent $A$ delegates to $B$ which subdelegates to $C$, the
chain of frames forms a directed acyclic graph of signed claims. Each
edge is independently verifiable. \zapaa{} does not require any single
party to be online during audit. This is in deliberate contrast to
the bearer-token model, where revocation of $B$'s token retroactively
invalidates the chain.

\subsection{Encrypted-to-recipient routing}

A frame is encrypted under the recipient's per-session AEAD key,
itself derived from the recipient's ML-KEM static key. A routing
gateway sees only the routing header (sender, recipient \rns{}
identifiers, frame size, sequence). The gateway cannot read or modify
content, removing an entire trusted compute base from the agent
fabric.

\section{Security}\label{sec:security}

We give formal statements here. Full proofs are in the companion
proof document~\cite{proof}.

\subsection{Definitions}

\begin{definition}[Agent non-repudiation]
For a frame $F$ produced in session $\mathit{sid}$ at sequence
$\mathit{seq}$, claimed sender $A$, payload $\hat{\rho}$, signature
$\sigma$: $F$ is non-repudiable if any third-party verifier $V$ holding
$\mathit{sig\_pk}_A$ can decide whether $\sigma$ is a valid ML-DSA-65
signature on $h(\hat{\rho}\,\Vert\,\mathit{ctx})$, with binding context
$\mathit{ctx}$ as defined above. The non-repudiation advantage of
$\mathcal{A}$ is the probability that $\mathcal{A}$ produces a triple
$(F^*, A^*, \hat{\rho}^*)$ accepted by $V$ as authored by $A^*$ when in
fact $A^*$ never signed it.
\end{definition}

\begin{definition}[Agent impersonation]
$\mathcal{A}$ impersonates an agent $A$ if $\mathcal{A}$ completes a
\zap{} handshake with an honest peer $T$ such that $T$ accepts the
session as authenticated to $A$ and $\mathcal{A}$ does not hold $A$'s
secret signing key.
\end{definition}

\begin{definition}[Cross-session replay]
A replay event occurs when $\mathcal{A}$ injects into an honest session
$\mathit{sid}_2$ a frame originally accepted in a distinct session
$\mathit{sid}_1$, and the recipient accepts it as fresh.
\end{definition}

\subsection{Theorems}

\begin{theorem}[Non-repudiation]\label{thm:nonrep}
For any PPT adversary $\mathcal{A}$ against the agent
non-repudiation game of \zapmcp{} or \zapaa{},
\[
\Adv^{\mathit{nonrep}}_{\zapmcp{}}(\mathcal{A}) \;\leq\;
q_s \cdot \Adv^{\mathit{eufcma}}_{\Sign,\mathit{ML\text{-}DSA\text{-}65}}(\mathcal{B})
+ q_h \cdot 2^{-256}
\]
where $q_s$ is the number of signing queries, $q_h$ the number of hash
queries, and $\mathcal{B}$ a constructed forger against ML-DSA-65 with
running time within an additive constant of $\mathcal{A}$.
\end{theorem}

\begin{theorem}[Agent impersonation infeasibility]\label{thm:imp}
For any PPT adversary $\mathcal{A}$ against \zapmcp{}/\zapaa{}
agent-impersonation,
\[
\Adv^{\mathit{imp}}_{\zap}(\mathcal{A}) \;\leq\;
\Adv^{\mathit{indcca}}_{\KEM,\mathit{ML\text{-}KEM\text{-}768}}(\mathcal{B}_1)
+ \Adv^{\mathit{eufcma}}_{\Sign,\mathit{ML\text{-}DSA\text{-}65}}(\mathcal{B}_2)
+ \frac{q_h^2}{2^{256}}.
\]
\end{theorem}

\begin{lemma}[Replay freshness]\label{lem:replay}
For session nonces of length $\nu$ bits and per-session frame counts
$\leq Q$,
\[
\Pr[\mathit{replay\ accepted}]
\;\leq\; \frac{Q(Q-1)}{2^{\nu+1}} + \mathit{negl}.
\]
With $\nu = 96$ and $Q \leq 2^{40}$ this is $\leq 2^{-17}$, well below
any operationally relevant threshold.
\end{lemma}

\begin{remark}
Theorems~\ref{thm:nonrep} and \ref{thm:imp} together imply that an
agent fabric built on \zapmcp{}/\zapaa{} delivers non-repudiable
provenance and bilateral authentication under standard post-quantum
assumptions. Lemma~\ref{lem:replay} forces $Q \cdot \nu$
budgeting in operational deployments; in practice this is satisfied
trivially by 96-bit nonces.
\end{remark}

\section{Wire Performance}\label{sec:perf}

\subsection{Baseline accounting}

Consider a representative \mcp{} \texttt{tools/call} request:
$\rho$ encodes to roughly 192 bytes of JSON. The HTTP+SSE+JWT baseline
adds the following overhead per request/response pair:

\begin{table}[h]
\centering
\small
\begin{tabular}{lr}
\toprule
\textbf{Layer} & \textbf{Bytes} \\
\midrule
HTTP request line                              & 24 \\
HTTP host header                               & 32 \\
HTTP \texttt{Authorization} (JWT, RS256)       & 920 \\
HTTP content-type / accept / user-agent / ...  & 180 \\
HTTP/2 HPACK savings (estimated)               & $-200$ \\
TLS record framing per record                  & 29 \\
SSE \texttt{event:}/\texttt{data:} framing    & 22 \\
HTTP/2 frame headers                           & 9 \\
\midrule
\textbf{HTTP+SSE+JWT total per call (one direction)} & \textbf{$\approx$1016} \\
\bottomrule
\end{tabular}
\end{table}

A second JWT travels with the response only when callbacks reauthenticate;
typical SSE keeps the original JWT context, so we do not count it on the
return path. We also do not count TCP/IP framing, which is identical on
both sides.

\subsection{\zap{} accounting}

A \zap{} frame carrying a signed \mcp{} JSON-RPC request consists of:

\begin{table}[h]
\centering
\small
\begin{tabular}{lr}
\toprule
\textbf{Field} & \textbf{Bytes} \\
\midrule
Frame magic + version                                  & 4 \\
Sequence number                                        & 8 \\
AEAD nonce (derived; on-wire nonce)                    & 12 \\
AEAD tag                                               & 16 \\
Length prefix                                          & 4 \\
Routing recipient ID (\rns{} hash)                     & 16 \\
ML-DSA-65 signature (in plaintext, AEAD-protected)     & 3309 \\
\midrule
\textbf{\zap{} per-frame fixed overhead}        & \textbf{60 (excl.\ signature)} \\
\textbf{\zap{} with per-frame ML-DSA-65 sig}    & \textbf{3369} \\
\bottomrule
\end{tabular}
\end{table}

The signature is dominated by ML-DSA-65's 3309-byte size. For
non-financial calls where third-party verifiability is not required, an
operator may select \emph{session-only} signatures (transcript signed
once at handshake, frames carry only AEAD tags); fixed per-frame
overhead is then 60 bytes versus HTTP+SSE+JWT's $\approx$1016 bytes ---
a 17$\times$ reduction. With per-frame ML-DSA-65 signatures the wire is
heavier than HTTP+JWT, but JWT carries no third-party verifiability:
the comparison is between a 3309-byte cryptographic record of intent
versus a 920-byte bearer assertion.

\subsection{Practical recommendation}

For \mcp{} tool calls of consequence (settlement, mutation, identity
issuance), per-frame ML-DSA-65 signatures are mandatory and the
cost is amortised over the audit value. For idempotent read-only
queries, session-signed frames are sufficient, and \zapmcp{} is on the
wire shorter than HTTP+SSE+JWT.

\begin{table}[h]
\centering
\small
\begin{tabular}{lcc}
\toprule
\textbf{Call type} & \textbf{HTTP+SSE+JWT (B)} & \textbf{\zapmcp{} (B)} \\
\midrule
read-only (e.g.\ \texttt{tools/list})  & 1016 & 60 \\
mutating, audited                       & 1016 & 3369 \\
\bottomrule
\end{tabular}
\end{table}

\section{Composition with \rns{}}

The full \zapmcp{} call sequence is:

\begin{enumerate}[leftmargin=2em,itemsep=0.25em]
\item agent looks up tool name $n$ in \rns{} (cached locally with TTL),
\item agent verifies record signature chain against the publisher root,
\item agent extracts $\mathit{kem\_pk}_T, \mathit{sig\_pk}_T,
\mathit{caps}_T$,
\item agent connects to $T$ at the network endpoint also published in
the record (an IP, hostname, or libp2p multiaddr),
\item agent runs the \zap{} 1-RTT handshake, receiving and verifying
$T$'s signed transcript,
\item agent issues \zapmcp{} frames as in \S\ref{sec:zap-mcp}.
\end{enumerate}

No HTTP layer participates. No JWT is ever issued. There is no
\emph{a priori} HTTPS endpoint. The identity of $T$ is its
$\mathit{kem\_pk}_T$, not any name resolved by DNS.

\section{Compliance Considerations}\label{sec:compliance}

\subsection{HIPAA Security Rule}

The HIPAA Security Rule~\cite{hipaa_security} requires
``transmission security'' (\S 164.312(e)) and ``audit controls'' (\S
164.312(b)). \zapmcp{} satisfies the encryption-in-transit requirement
under hybrid ML-KEM-768 + X25519, and audit controls map to the per-frame
signed log: a clinical agent's request to read a record is a third-party
verifiable artifact, retainable for the statutory six years.

\subsection{SOC~2}

SOC~2~\cite{soc2} Common Criteria CC6.7 (transmission of information)
and CC7.2 (system monitoring) are addressed identically. The
distinguishing benefit is that \zap{}-logged events are
\emph{independently} verifiable; SOC~2 auditors can replay the
verification without trusting the operator's logging pipeline.

\subsection{GDPR Article 30}

Article~30~\cite{gdpr_art30} requires ``records of processing
activities'' kept by the controller and processor. A \zap{}-wrapped
agent fabric naturally produces such records: each tool call's request
and response are signed, dated (via session-bound monotonic sequence),
and bound to the agent and tool identity. This is materially stronger
than the typical Article~30 ``description of categories of processing''
prose, which is non-cryptographic.

\subsection{Long-term retention and the quantum horizon}

Medical, legal, and financial records routinely have retention periods
of 6--75 years. Any TLS-only protection of those records today must be
modelled under harvest-now-decrypt-later~\cite{harvest_now}: an attacker
who records the TLS handshake in 2026 and obtains a CRQC by, say, 2035,
recovers all session content. Hybrid ML-KEM-768 keying eliminates this
class of risk for any traffic encrypted under \zap{}.

\section{Deployment Notes}

\subsection{Existing \mcp{} servers}

A backwards-compatible adapter pattern is straightforward: a thin
\zapmcp{} sidecar terminates frames, validates the agent identity and
capability claims, then forwards plain JSON-RPC over a UNIX domain
socket to the existing \mcp{} server process. The sidecar adds
$\approx$200 lines and runs in process; this is the recommended
migration path for tool authors.

\subsection{Existing \aatwo{} agents}

\aatwo{} agents accept agent cards by URL. \zapaa{} extends this with
\rns{}-resolved cards. Agents that wish to interoperate with both can
publish a \aatwo{}-card-by-URL whose contents are byte-for-byte the
canonical CBOR encoding of the \rns{} record content; verifiers can
then cross-check.

\subsection{Key management}

Long-term ML-DSA-65 signing keys are held in HSMs or platform secure
enclaves. Per-session handshake keys are derived from the long-term keys
plus ephemeral KEM material; session keys are zeroised at session
close. Backup of long-term keys is by Shamir secret sharing across
operator-distinct custodians; in regulated environments a regulator
may hold one share with auditable retrieval policy.

\section{Gap Items and Future Work}\label{sec:gaps}

We are explicit about what \zapmcp{}/\zapaa{} does not yet solve.

\begin{enumerate}[leftmargin=2em,itemsep=0.25em]
\item \textbf{Revocation distribution.} If a long-term agent signing
key is compromised, every cached \rns{} record pointing at it must be
poisoned. We currently rely on short \rns{} TTLs (default 600 seconds)
plus a transparency-log-style revocation feed. A formal CRL-equivalent
distribution scheme suitable for a fleet of $10^6$ agents is open work.
The trade-off space is short TTLs (chatty resolvers) versus
signed delta updates over a gossip overlay (complex, risk of
gossip-poisoning).
\item \textbf{Capability claim semantics.} \rns{} records carry a
capability vector, but the schema is application-specific. A canonical
ontology for tool capabilities (read-only, idempotent, mutating,
financial-impact) would let policy be expressed transport-side rather
than per-application.
\item \textbf{Indirect prompt injection.} \zapmcp{} attributes
tool output to a tool, but does not detect when that output is a
prompt-injection payload. Mitigation belongs at the model integration
layer; \zap{} provides forensic attribution.
\item \textbf{Quantum-safe \rns{} root.} The \rns{} publisher root key
itself must be ML-DSA-65 (or SLH-DSA-256s for ultra-long horizon). A
governance procedure for rotating that root in a federated agent fabric
is open.
\item \textbf{Tamarin/ProVerif analysis.} We have proven properties in
the computational model (Section~\ref{sec:security}). A symbolic model
verification along the lines of \cite{cremers2017ake} for full
\zapmcp{} session establishment is in progress.
\item \textbf{Side-channel resistance of agent-side ML-DSA-65.} On
phones and laptops running an LLM agent, ML-DSA-65 signing must run in
constant time; the reference implementation does, but operator-grade
deployments should be verified against the target microarchitecture.
\end{enumerate}

\section{Related Work}

\paragraph{TLS-style mutual authentication.} TLS~\cite{rfc8446} with
client certificates achieves mutual authentication but is bound to
hostnames and X.509 PKIs~\cite{rfc5280}; agent populations are
fundamentally not hostname-shaped. ACME~\cite{rfc8555} and short-lived
certificates partially address dynamism but inherit the X.509 trust
model.

\paragraph{DIDs and Verifiable Credentials.} W3C
DIDs~\cite{didcore,didweb,didkey} and Verifiable
Credentials~\cite{vc20} provide a self-sovereign identity model close
in spirit to \rns{}. \rns{} differs by enforcing that resolution
yields a KEM public key and binds transport to it; this collapses the
``identity proof'' and ``transport key'' into a single artifact and
removes a class of failure modes where DID resolution and transport
authentication disagree.

\paragraph{COSE/JOSE signed envelopes.} COSE~\cite{cose} and
JWE~\cite{jose} provide signed and encrypted message envelopes. They
do not specify a session establishment, do not pin transport identity
to KEM keys, and do not provide replay freshness binding.
\zap{} uses COSE-style canonicalisation under the covers but adds the
session and identity layers.

\paragraph{Agent governance literature.} \cite{agent_audit_2024}
articulates non-cryptographic governance practices for agentic AI.
Cryptographic audit substrate as offered by \zapmcp{}/\zapaa{} is
complementary and provides the substrate on which governance practices
can be enforced.

\section{Conclusion}

The agent fabric is the new universal joint of computation. Its
security cannot be the security of HTTP plus a bearer token. We have
shown that wrapping \mcp{} and \aatwo{} in \zap{} delivers
post-quantum confidentiality, mutual authentication of agents and
tools by KEM identity, non-repudiation of every message via per-frame
ML-DSA-65 signatures, and provenance suitable for HIPAA, SOC~2, and
GDPR. The construction is shorter on the wire than the HTTP+SSE+JWT
baseline for read-only calls and competitive for audit-grade mutating
calls. Open work centres on revocation distribution at fleet scale,
canonical capability ontologies, and symbolic model verification of
the full handshake. We invite implementors of agent platforms to treat
this work as a starting specification.

\bibliographystyle{plain}
\bibliography{refs}

\end{document}
